\begin{sheetframe}[50]{violet}{04｜距離空間・位相}{定義・保存・短い証明}

\begin{columns}[T,onlytextwidth]
\begin{column}{.285\textwidth}
\subhead{距離空間 $(X,d)$}
\[
d:X\times X\longrightarrow[0,\infty),\qquad
x,y,z\in X
\]
\[
\begin{aligned}
d(x,y)=0&\iff x=y,\\
d(x,y)&=d(y,x),\\
d(x,z)&\le d(x,y)+d(y,z).
\end{aligned}
\]
\[
\boxed{B(x;r):=\{y\in X:d(x,y)<r\}\quad(r>0)}
\]
\end{column}

\begin{column}{.335\textwidth}
\subhead{位相空間 $(X,\mathcal T_X)$}
\[
2^X:=\{A:A\subset X\},\qquad
\underbrace{\mathcal T_X}_{X\text{ 上の位相}}\subset 2^X
\]
\[
\begin{aligned}
&\varnothing,X\in\mathcal T_X,\\
&(U_\alpha)_{\alpha\in I}\subset\mathcal T_X
 \Rightarrow \bigcup_{\alpha\in I}U_\alpha\in\mathcal T_X,\\
&U,V\in\mathcal T_X
 \Rightarrow U\cap V\in\mathcal T_X.
\end{aligned}
\]
ここで $I$ は任意の添字集合。
この3条件を満たす $\mathcal T_X$ を位相、$(X,\mathcal T_X)$ を位相空間、$\mathcal T_X$ の要素を開集合という。
\end{column}

\begin{column}{.35\textwidth}
\subhead{距離が作る位相・連続写像}
\[
\mathcal T_d
:=\left\{U\subset X:
\begin{array}{l}
\forall x\in U,\ \exists r>0,\\[-.5mm]
B(x;r)\subset U
\end{array}\right\}.
\]
距離空間 $(X,d)$ は $(X,\mathcal T_d)$ という位相空間になる。
\begin{center}
\begin{tikzpicture}[x=.43mm,y=.43mm]
  \draw[accent,line width=.55pt] (3,0) ellipse (16 and 6);
  \draw[accent!75,line width=.45pt] (3,0) ellipse (7 and 3);
  \node[tinylabel] at (-10,4) {$X$};
  \node[tinylabel] at (3,0) {$f^{-1}(U)$};
  \draw[accent,line width=.55pt] (52,0) ellipse (16 and 6);
  \draw[accent!75,line width=.45pt] (52,0) ellipse (7 and 3);
  \node[tinylabel] at (39,4) {$Y$};
  \node[tinylabel] at (52,0) {$U$};
  \draw[flow] (20,1)--node[above,tinylabel] {$f$}(35,1);
\end{tikzpicture}
\end{center}
\[
\boxed{
\begin{gathered}
f:(X,\mathcal T_X)\to(Y,\mathcal T_Y),\\
f\text{ 連続}
\iff
\forall U\in\mathcal T_Y,\ f^{-1}(U)\in\mathcal T_X.
\end{gathered}}
\]
同値な見方：$Y$ の閉集合 $F$ に対し、$f^{-1}(F)$ は $X$ で閉。
\end{column}
\end{columns}

\vspace{.3mm}{\color{line}\rule{\textwidth}{.35pt}}\vspace{.25mm}

\begin{columns}[T,onlytextwidth]
\begin{column}{.315\textwidth}
\subhead{開・閉・内部・閉包・境界}
\begin{center}
\begin{tikzpicture}[x=.62mm,y=.62mm]
  \draw[muted,line width=.4pt] (-19,-8) rectangle (20,9);
  \node[tinylabel] at (-16,6) {$X$};
  \fill[accent!9] plot[smooth cycle,tension=.75]
    coordinates {(-12,-4)(-6,-6)(4,-4)(8,2)(1,6)(-10,4)};
  \draw[accent,line width=.6pt] plot[smooth cycle,tension=.75]
    coordinates {(-12,-4)(-6,-6)(4,-4)(8,2)(1,6)(-10,4)};
  \node[labelnode] at (-3,0) {$A$};
  \fill[coral] (11,2) circle (.8) node[above right,tinylabel] {$x$};
  \draw[coral,dashed,line width=.45pt] (11,2) circle (5);
  \node[tinylabel,anchor=west] at (16,2) {$B(x;r)\cap A\ne\varnothing$};
\end{tikzpicture}
\end{center}
\[
\begin{aligned}
U\subset X\text{ が開}
&\iff U\in\mathcal T_X,\\
F\subset X\text{ が閉}
&\iff X\setminus F\in\mathcal T_X.
\end{aligned}
\]
\[
\begin{aligned}
x\in X,\quad A\subset X:\qquad
x\in A^\circ
&\iff \exists U\in\mathcal T_X:
     x\in U\subset A,\\
x\in\overline A
&\iff \forall U\in\mathcal T_X,\
     x\in U\Rightarrow U\cap A\ne\varnothing.
\end{aligned}
\]
\[
\partial A=\overline A\setminus A^\circ,\qquad
\overline A=\bigcap_{\substack{A\subset F\subset X\\F\text{ は }X\text{ で閉}}}F.
\]
位相空間での収束は
\[
x_n\to x
\iff
\forall U\in\mathcal T_X,\quad
x\in U\Rightarrow
\exists N\in\N,\ \forall n\ge N,\ x_n\in U.
\]
距離空間では
\[
x\in\overline A
\iff
\forall r>0,\ B(x;r)\cap A\ne\varnothing
\iff
\exists a_n\in A,\ a_n\to x.
\]

\separator
\subhead{相対位相・商位相}
部分集合 $S\subset X$ に対して
\[
\boxed{
\underbrace{\mathcal T_S}_{S\text{ 上の相対位相}}
:=\{S\cap U:U\in\mathcal T_X\}\subset 2^S.}
\]
たとえば $X=\R,\ S=[0,1]$ なら
\[
[0,\tfrac12)=S\cap(-1,\tfrac12)\in\mathcal T_S.
\]
これは $S$ では開だが、$\R$ では開でない。

$\sim$ を $X$ 上の同値関係とし、$Q=X/{\sim}$、$\pi:X\to Q$, $\pi(x)=[x]$ とする。
\[
\boxed{
\underbrace{\mathcal T_Q}_{Q\text{ 上の商位相}}
:=\{V\subset Q:\pi^{-1}(V)\in\mathcal T_X\}.}
\]

\separator
\subhead{ハウスドルフ}
\[
\forall x,y\in X,\quad x\ne y\Rightarrow
\exists U,V\in\mathcal T_X:
\quad
\begin{cases}
x\in U,\ y\in V,\\
U\cap V=\varnothing.
\end{cases}
\]
距離空間はハウスドルフ。したがって極限は一意。
\end{column}

\begin{column}{.37\textwidth}
\subhead{コンパクトの定義}
\begin{center}
\begin{tikzpicture}[x=.58mm,y=.58mm]
  \fill[accent!8] plot[smooth cycle,tension=.75]
    coordinates {(-22,-5)(-13,-10)(1,-8)(18,-2)(15,7)(2,10)(-16,7)};
  \draw[accent,line width=.65pt] plot[smooth cycle,tension=.75]
    coordinates {(-22,-5)(-13,-10)(1,-8)(18,-2)(15,7)(2,10)(-16,7)};
  \foreach \x/\y in {-16/1,-7/6,2/4,11/4,14/-3,4/-6,-8/-5}{
    \draw[muted!65,line width=.3pt] (\x,\y) circle (7);
  }
  \foreach \x/\y in {-16/1,2/4,14/-3,4/-6}{
    \draw[coral,line width=.7pt] (\x,\y) circle (7);
  }
  \node[tinylabel] at (0,-13) {赤い4個で覆える};
\end{tikzpicture}
\end{center}
\[
\begin{gathered}
K\subset X,\quad U_\alpha\in\mathcal T_X,\\
K\subset\bigcup_{\alpha\in I}U_\alpha
\ \Longrightarrow\
\exists\alpha_1,\ldots,\alpha_m:\
K\subset\bigcup_{i=1}^mU_{\alpha_i}.
\end{gathered}
\]
\[
\begin{aligned}
\text{距離空間：}\quad
K\text{ コンパクト}
&\iff
\text{任意の }(x_n)\subset K\\[-.4mm]
&\quad\text{が }K\text{ 内に収束する部分列を持つ},\\
K\subset\R^n:\quad
K\text{ コンパクト}
&\iff K\text{ は閉かつ有界}.
\end{aligned}
\]

\separator
\subhead{性質と短い証明}
\minorhead{$K$ がコンパクトで、$F\subset K$ が $K$ で閉なら $F$ はコンパクト}
\[
\begin{aligned}
U_\alpha\in\mathcal T_K,\quad
F\subset\bigcup_\alpha U_\alpha
&\Rightarrow
\{U_\alpha\}\cup\{K\setminus F\}
   \text{ が }K\text{ を覆う}\\
&\Rightarrow
\text{有限個で }F\text{ を覆える}.
\end{aligned}
\]
\minorhead{$f:K\to Y$ が連続なら、コンパクト $K$ の像 $f(K)$ もコンパクト}
\[
\begin{aligned}
V_\alpha\in\mathcal T_Y,\quad
f(K)\subset\bigcup_\alpha V_\alpha
&\Rightarrow
K\subset\bigcup_\alpha f^{-1}(V_\alpha)\\
&\Rightarrow
\text{有限個の }V_\alpha\text{ で }f(K)\text{ を覆える}.
\end{aligned}
\]
\minorhead{$X$ がハウスドルフなら、コンパクト $K\subset X$ は閉}
\[
\begin{aligned}
y\notin K,\ x\in K
&\Rightarrow
\exists U_x,V_x\in\mathcal T_X:\quad
x\in U_x,\ y\in V_x,\ U_x\cap V_x=\varnothing,\\
K\subset\bigcup_{x\in K}U_x
&\Rightarrow
K\subset U_{x_1}\cup\cdots\cup U_{x_m},\\
V=V_{x_1}\cap\cdots\cap V_{x_m}
&\Rightarrow
y\in V\subset X\setminus K.
\end{aligned}
\]
よって $X\setminus K$ は開、したがって $K$ は閉。
\end{column}

\begin{column}{.295\textwidth}
\subhead{連結・道連結の定義}
\begin{center}
\begin{tikzpicture}[x=.55mm,y=.55mm]
  \fill[accent!9] (-18,0) circle (7); \draw[accent] (-18,0) circle (7);
  \fill[accent!9] (0,0) circle (7); \draw[accent] (0,0) circle (7);
  \node[tinylabel] at (-9,-10) {$X=U\sqcup V$：非連結};
  \fill[accent!9] plot[smooth cycle,tension=.8]
    coordinates {(14,-5)(21,-8)(31,-5)(32,5)(22,8)(12,4)};
  \draw[accent] plot[smooth cycle,tension=.8]
    coordinates {(14,-5)(21,-8)(31,-5)(32,5)(22,8)(12,4)};
  \draw[coral,line width=.7pt] (16,-2)..controls(21,6)and(27,-5)..(30,3);
  \node[tinylabel] at (23,-10) {点を道で結ぶ};
\end{tikzpicture}
\end{center}
\[
\begin{aligned}
X\text{ が非連結}
\iff
\exists U,V\in\mathcal T_X:\quad
&U\ne\varnothing,\ V\ne\varnothing,\\
&U\cap V=\varnothing,\quad X=U\cup V.
\end{aligned}
\]
$X$ が連結とは、上の $U,V$ が存在しないこと。
\[
\begin{aligned}
X\text{ が道連結}
\iff
\forall x,y\in X,\ \exists\gamma:[0,1]\to X:\quad
&\gamma\text{ は連続},\\
&\gamma(0)=x,\quad\gamma(1)=y.
\end{aligned}
\]
\minorhead{道連結 $\Rightarrow$ 連結}
分離 $X=U\sqcup V$ があれば、$x\in U,\ y\in V$ を結ぶ $\gamma$ により
$[0,1]=\gamma^{-1}(U)\sqcup\gamma^{-1}(V)$ となり、区間の連結性に反する。
\minorhead{連続像は連結性を保つ}
$f:X\to Y$ が連続で $f(X)=U\sqcup V$ と分離できるなら、
$X=f^{-1}(U)\sqcup f^{-1}(V)$ となるので、$X$ の連結性に反する。

\separator
\subhead{集合までの距離}
$(X,d)$ を距離空間、$\varnothing\ne A\subset X$ とする。
\[
\begin{aligned}
d(x,A)&=\inf_{a\in A}d(x,a),\\
|d(x,A)-d(y,A)|&\le d(x,y),\\
d(x,A)=0&\iff x\in\overline A.
\end{aligned}
\]
$A$ がコンパクトなら、連続関数 $a\mapsto d(x,a)$ は最小値を取るので、
$d(x,A)=d(x,a_0)$ となる最近点 $a_0\in A$ が存在する。

\separator
\subhead{よく使う帰結と反例}
\[
f\text{ が同相写像}
\iff
f\text{ は全単射で }f,\ f^{-1}\text{ がともに連続}.
\]
\[
\begin{gathered}
K\text{ コンパクト},\ Y\text{ ハウスドルフ},\\
f:K\to Y\text{ 連続全単射}
\ \Longrightarrow\ f\text{ は同相写像}.
\end{gathered}
\]
理由：閉集合 $F\subset K$ はコンパクトなので、$f(F)$ はコンパクト、したがって $Y$ で閉。よって $f^{-1}$ は連続。
\[
\begin{aligned}
f:\R\to\R,\quad f(x)=x^2:\qquad
&f((-1,1))=[0,1)
&&\text{は開でない},\\
f:\R\to\{0\}:\qquad
&f^{-1}(\{0\})=\R
&&\text{はコンパクトでない}.
\end{aligned}
\]
\end{column}
\end{columns}
\end{sheetframe}
